\section{重新理解决策窗口}

这一生物学层面使决策窗口的含义更加清晰。
决策窗口不只是一个关于动机的口号。
它可以被理解为一个短暂区间：在这一时段内，当前机制已经削弱到足以让一个小干预改变轨迹，
而旧的吸引盆尚未重新夺回控制。

\begin{example}[把入睡视为一种转变模型]
这篇关于睡眠转变的文献之所以有用，恰恰在于它并不是手稿自行发明的隐喻。
被索引的 Nature Neuroscience 文章把入睡处理为一种可预测的分岔动力学。
这并不意味着每一种行为切换都在字面上使用同样的机制，
但它确实使我们在日常生活中可以提出一个更尖锐的问题：
究竟是哪一个控制参数正在被推动，当前机制距离失去稳定性还有多近，
以及在路径真正弯折之前，需要多大幅度的干预？
\end{example}